\documentclass[12pt]{article}
\usepackage{amsmath, amssymb}

\begin{document}

\title{Lecture Scribe: Inequalities and Tail Bounds}
\author{Course: CSE400 – Fundamentals of Probability in Computing \\ Instructor: Dhaval Patel \\ Reference Text: Introduction to Probability for Computing}
\date{}
\maketitle

\section*{1. Introduction and Outline}

The lecture covers the following topics in sequence:

\begin{itemize}
\item Tail and motivation of tail bounds
\item Markov’s Inequality (definition, proof, example)
\item Chebyshev’s Inequality (definition, proof, example)
\item Chernoff Bound and Poisson tail (definition, proof, example)
\item Jensen’s Inequality and convex functions
\item Case study and system-level understanding
\end{itemize}

\section*{2. Tail of a Random Variable}

\textbf{Definition}

The tail of a random variable $X$ is defined as:

\[
P(X > x)
\]

This represents the probability that the random variable takes values larger than a threshold $x$.

\section*{3. Motivation: Why Study Tails?}

The lecture motivates tail probabilities through real-world system events:

\begin{itemize}
\item Jobs delayed more than 24 hours
\item Packets dropped due to full buffer
\end{itemize}

\textbf{Key Idea}

Mean $\rightarrow \rightarrow$ captures average behavior (easy to compute) \\
Tail $\rightarrow \rightarrow$ captures rare extreme events (hard to compute)

\textbf{Why Are Tails Hard to Compute?}

Requires summing many small probabilities

\section*{4. Tail Example 1: Jobs Distributed Across Servers}

\textbf{Problem Statement}

Suppose $n$ jobs are distributed among $n$ servers at random. \\
Find:

\[
P(\text{a particular server gets } \geq k \text{ jobs})
\]

\textbf{Intuition}

\begin{itemize}
\item Jobs are assigned randomly
\item Most servers get few jobs
\item Occasionally, one server gets many jobs
\end{itemize}

This corresponds to a tail event (rare overload).

\textbf{Model}

\[
X \sim \text{Binomial}(n,p)
\]

\textbf{Exact Probability}

\[
P(X \geq k) = \sum_{i=k}^{n} \binom{n}{i} p^i (1-p)^{n-i}
\]

\section*{5. Tail Example 2: Poisson Arrivals}

\textbf{Problem Statement}

Jobs arrive to a datacenter according to a Poisson process with rate $\lambda$ jobs/hour. \\
Find:

\[
P(\text{at least } k \text{ jobs arrive in one hour})
\]

\textbf{Intuition}

\begin{itemize}
\item Jobs arrive randomly over time
\item Most hours have moderate arrivals
\item Occasionally, bursts occur
\end{itemize}

This corresponds to a tail event (sudden spike in load).

\textbf{Model}

\[
X \sim \text{Poisson}(\lambda)
\]

\textbf{Exact Probability}

\[
P(X \geq k) = \sum_{i=k}^{\infty} e^{-\lambda} \frac{\lambda^i}{i!}
\]

\section*{6. Why Tail Bounds Are Needed}

\textbf{Observation}

Exact tail probabilities are often difficult to compute because:

\begin{itemize}
\item They involve long summations
\item Expressions are computationally expensive
\end{itemize}

\textbf{Examples of Exact Computations}

Poisson:
\[
P(X \geq k) = \sum_{i=k}^{\infty} e^{-\lambda} \frac{\lambda^i}{i!}
\]

Binomial:
\[
P(X \geq k) = \sum_{i=k}^{n} \binom{n}{i} p^i (1-p)^{n-i}
\]

These are complicated, long sums.

\section*{7. Key Idea: Tail Bounds}

Instead of computing exact values:

\begin{itemize}
\item We compute an upper bound
\item Easier to evaluate
\item Still provides a guarantee
\end{itemize}

\textbf{Formulation}

Exact: 
\[
P(X \geq k) = ?
\]

Bound:
\[
P(X \geq k) \leq \text{simple expression}
\]

\textbf{Key Insight}

Approximate safely instead of computing exactly.

\section*{8. Running Example (Used Throughout Lecture)}

\textbf{Goal}

Develop progressively tighter tail bounds and compare them on the same example.

\textbf{Experiment}

Flip a fair coin $n$ times

\textbf{Random Variable}

\[
X = \text{number of heads}
\]

\[
X \sim \text{Binomial}\left(n, \frac{1}{2}\right)
\]

\textbf{Mean}

\[
E[X] = \frac{n}{2}
\]

\textbf{Tail Question}

\[
P\left(X \geq \frac{3n}{4}\right)
\]

\textbf{Interpretation}

This corresponds to the tail region: \\
``getting unusually many heads.''

\section*{9. Markov’s Inequality}

\textbf{Key Idea}

\begin{itemize}
\item Uses only the mean
\item Large values are rare
\end{itemize}

\textbf{Example Motivation}

If average = 10, then very large values (e.g., 100+) are rare.

\textbf{Statement (Markov’s Inequality)}

For a non-negative random variable $X$:

\[
P(X \geq a) \leq \frac{E[X]}{a}
\]

\section*{10. Summary of Concepts Covered}

Tail probabilities measure rare events:

\[
P(X > x)
\]

Exact computation is difficult due to long summations.

Tail bounds provide:

\begin{itemize}
\item Simpler computation
\item Guaranteed upper bounds
\end{itemize}

Running example:

\[
X \sim \text{Binomial}(n, 1/2)
\]

Study:

\[
P(X \geq 3n/4)
\]

First bound introduced:

Markov’s Inequality

\end{document}